\section{Квадратичная аппроксимация}

В этом разделе строится полином наилучшего среднеквадратического отклонения
степени $m = 2$, то есть квадратичная аппроксимация
$P_2(x) = a_0 + a_1 x + a_2 x^2$.
Нормальная система МНК для её коэффициентов:
\[
\begin{cases}
a_0 n + a_1 \displaystyle\sum x_i + a_2 \sum x_i^2
    = \sum y_i, \\[6pt]
a_0 \displaystyle\sum x_i + a_1 \sum x_i^2 + a_2 \sum x_i^3
    = \sum y_i x_i, \\[6pt]
a_0 \displaystyle\sum x_i^2 + a_1 \sum x_i^3 + a_2 \sum x_i^4
    = \sum y_i x_i^2.
\end{cases}
\]

\subsection{Дополнительные суммы}

Помимо сумм, найденных в разделе~2, потребуются ещё $\sum x_i^3$, $\sum x_i^4$ и
$\sum y_i x_i^2$. Они вычислены в таблице~\ref{tab:aux} и равны:
\[
\sum x_i^3 = 44{,}1000, \quad \sum x_i^4 = 72{,}2666, \quad
\sum y_i x_i^2 = 52{,}5719.
\]

\subsection{Решение нормальной системы}

Подставляя все суммы, нормальная система принимает вид:
\[
\begin{cases}
20\,a_0 + 21{,}0000\,a_1 + 28{,}7000\,a_2 = 47{,}2100, \\
21{,}0000\,a_0 + 28{,}7000\,a_1 + 44{,}1000\,a_2 = 42{,}7030, \\
28{,}7000\,a_0 + 44{,}1000\,a_1 + 72{,}2666\,a_2 = 52{,}5719.
\end{cases}
\]

На листинге~\ref{lst:quadratic} приведён код для решения этой системы.

\begin{lstlisting}[caption={Квадратичная аппроксимация МНК}, label=lst:quadratic]
import numpy as np

x = np.array([0.1,0.2,0.3,0.4,0.5,0.6,0.7,0.8,0.9,1.0,
              1.1,1.2,1.3,1.4,1.5,1.6,1.7,1.8,1.9,2.0])
y = np.array([3.15,3.04,3.02,2.97,2.87,2.98,2.81,2.70,2.66,2.50,
              2.60,2.36,2.09,2.07,2.01,1.81,1.53,1.64,1.29,1.11])

n    = len(x)
sx   = np.sum(x)
sx2  = np.sum(x**2)
sx3  = np.sum(x**3)
sx4  = np.sum(x**4)
sy   = np.sum(y)
sxy  = np.sum(x * y)
sx2y = np.sum(x**2 * y)

# Нормальная система: A * a = b
A = np.array([[n,   sx,  sx2],
              [sx,  sx2, sx3],
              [sx2, sx3, sx4]])
b = np.array([sy, sxy, sx2y])

a0, a1, a2 = np.linalg.solve(A, b)

# Аппроксимирующий полином и отклонение
P2     = a0 + a1 * x + a2 * x**2
delta2 = np.sqrt(np.sum((y - P2)**2) / n)

print(f"a0 = {a0:.4f},  a1 = {a1:.4f},  a2 = {a2:.4f}")
print(f"P2(x) = {a0:.4f} + ({a1:.4f})*x + ({a2:.4f})*x^2")
print(f"delta2(P2) = {delta2:.4f}")
\end{lstlisting}

Решение системы даёт коэффициенты:
\[
a_0 = 3{,}1147, \quad a_1 = -0{,}1323, \quad a_2 = -0{,}4287.
\]
Квадратичная аппроксимирующая функция:
\[
P_2(x) = 3{,}1147 - 0{,}1323\,x - 0{,}4287\,x^2.
\]

\subsection{Вычисление наилучшего среднеквадратического отклонения}

Наилучшее среднеквадратическое отклонение:
\[
\delta_2(P_2) = \sqrt{\frac{1}{n}\sum_{i=1}^{n}(y_i - P_2(x_i))^2}.
\]
Сумма квадратов отклонений $\sum(y_i - P_2(x_i))^2 = 0{,}1059$, поэтому:
\[
\delta_2(P_2) = \sqrt{\frac{1}{20} \cdot 0{,}1059} = 0{,}0728.
\]

Таблица~\ref{tab:dev} содержит значения $P_1(x_i)$, $P_2(x_i)$ и квадраты
отклонений для каждого узла.

\begin{table}[h!]
\centering
\small
\begin{tabular}{|c|c|c|c|c|c|c|}
\hline
$i$ & $x_i$ & $y_i$ & $P_1(x_i)$ & $(y_i - P_1)^2$ & $P_2(x_i)$ & $(y_i - P_2)^2$ \\
\hline
1  & 0,1 & 3,1500 & 3,3416 & 0,0367 & 3,0972 & 0,0028 \\
2  & 0,2 & 3,0400 & 3,2383 & 0,0393 & 3,0711 & 0,0010 \\
3  & 0,3 & 3,0200 & 3,1350 & 0,0132 & 3,0364 & 0,0003 \\
4  & 0,4 & 2,9700 & 3,0318 & 0,0038 & 2,9932 & 0,0005 \\
5  & 0,5 & 2,8700 & 2,9285 & 0,0034 & 2,9414 & 0,0051 \\
6  & 0,6 & 2,9800 & 2,8252 & 0,0240 & 2,8810 & 0,0098 \\
7  & 0,7 & 2,8100 & 2,7219 & 0,0078 & 2,8120 & 0,0000 \\
8  & 0,8 & 2,7000 & 2,6187 & 0,0066 & 2,7344 & 0,0012 \\
9  & 0,9 & 2,6600 & 2,5154 & 0,0209 & 2,6483 & 0,0001 \\
10 & 1,0 & 2,5000 & 2,4121 & 0,0077 & 2,5536 & 0,0029 \\
11 & 1,1 & 2,6000 & 2,3089 & 0,0848 & 2,4503 & 0,0224 \\
12 & 1,2 & 2,3600 & 2,2056 & 0,0238 & 2,3385 & 0,0005 \\
13 & 1,3 & 2,0900 & 2,1023 & 0,0002 & 2,2181 & 0,0164 \\
14 & 1,4 & 2,0700 & 1,9991 & 0,0050 & 2,0891 & 0,0004 \\
15 & 1,5 & 2,0100 & 1,8958 & 0,0130 & 1,9515 & 0,0034 \\
16 & 1,6 & 1,8100 & 1,7925 & 0,0003 & 1,8054 & 0,0000 \\
17 & 1,7 & 1,5300 & 1,6892 & 0,0254 & 1,6507 & 0,0146 \\
18 & 1,8 & 1,6400 & 1,5860 & 0,0029 & 1,4874 & 0,0233 \\
19 & 1,9 & 1,2900 & 1,4827 & 0,0371 & 1,3155 & 0,0006 \\
20 & 2,0 & 1,1100 & 1,3794 & 0,0726 & 1,1350 & 0,0006 \\
\hline
$\sum$ & & & & 0,4286 & & 0,1059 \\
\hline
\end{tabular}
\caption{Значения аппроксимирующих полиномов и квадраты отклонений.}
\label{tab:dev}
\end{table}